\section{Introduction}
\label{sec:introduction}

Binary encoding is a major cost in correlation generators and code-based
proof systems. Pseudorandom correlation generators (PCGs) apply public
generator matrices to expand sparse correlations into OT and VOLE
correlations~\cite{bcgi18,bcg19pcg}. Code-based proof systems encode rows
of large arrays, with distance supporting soundness~\cite{brakedown23,blaze25}.
These applications need constant rate and linear distance, but also fast
encoding, efficient transposed encoding, and regular memory access.

Random linear codes provide strong distance but require dense computation
under direct encoding. Structure can reduce arithmetic and data movement;
it also introduces dependencies that complicate distance proofs.
We introduce \emph{Single-Permutation INterleaved} (SPIN) codes to combine
quantitative distance guarantees with efficient evaluation in both directions.

For an admissible output length $N$, a SPIN encoder has the form
\begin{equation}
  E_N := I_N\circ\Pi_N\circ O_N.
  \label{eq:intro-spin-composition}
\end{equation}
The outer $O_N$ maps the message to $N$ intermediate bits, $\Pi_N$ permutes
these coordinates, and the recursive linear inner $I_N$ preserves length.
The outer gives sparse messages an initial weight floor; the interleaver
spreads the resulting support before recursive encoding. The phrase \emph{single permutation} describes this
interface: block shuffles, a transpose, and region shuffles may implement
$\Pi_N$, but their composition is one global interleaver.

\subsection{Results and Approach}

\paragraph{Distance and Encoding Work.}
Every realization of our main construction, Structured SPIN, has rate $1/2$.
Its minimum distance exceeds $0.11N$ with probability $1-o(1)$ over setup.
Ordinary and transposed encoding require $O(N)$ bit operations at the family's
native lengths. The distance and rate claims are formally verified in Lean,
including the construction's random setup, the analytic estimates, and the
required numerical certificates. Appendix~\ref{app:structured-lean} describes
the formalization.
For finite lengths, we choose constituents for concrete performance
rather than require the asymptotic outer. A fixed BCH-derived
$[256,128,\ge38]$ constituent gives relative distance greater than $0.10$
with setup-failure probability below $2^{-40}$ at message lengths
$2^{16},2^{18},2^{20},2^{22},2^{24}$.
The certificate uses rigorous spectrum inequalities, not an assumed exact
BCH spectrum. Section~\ref{sec:finite-quarter} also gives two certified
rate-$1/4$ operating points.

\paragraph{Analysis and Structure.}
For this class of encoder, sparse messages typically cause low minimum
distance. The encoder must spread their nonzero bits widely without dense
matrix multiplication.
First, the outer encodes short message blocks independently, so even a block
with a single nonzero bit produces several nonzero bits.
Next, the permutation spreads the bits of each encoded block across distinct
regions.
Finally, the inner processes this stream with a small running state.
Cheap XOR updates let an input bit influence many later outputs, spreading
local changes over a long codeword.
Our proof establishes the distance guarantees above by
bounding the probability that any nonzero message produces too few nonzero
output bits. Our contribution is this efficient combination and its distance
analysis; its ingredients build on established ideas from concatenated codes.

Section~\ref{sec:spin-intuition} develops the idea through two simpler
supporting constructions: Accumulator SPIN uses a running sum as its inner, and Random SPIN
illustrates what stronger random mixing can achieve. Structured SPIN
matches Random SPIN's $11\%$ distance guarantee while reducing the direct
encoding work from $\Theta(N\log N)$ to $O(N)$.
The complete proofs for these two supporting constructions are in
Appendix~\ref{app:supporting-constructions}.
The minimal accumulator instance also answers the logarithmic-outer-memory
question of Bazzi, Mahdian, and Spielman~\cite{bms09} in the
terminated-encoder interpretation of Appendix~\ref{sec:bms-question}.

\paragraph{Implementation and Applications.}
On one Ryzen 9 7950X thread, SPIN's transposed encoder computes
$2^{20}$ $128$-bit correlation blocks in about $9.3$\,ms after precomputation,
approximately $3.4\times$ faster than the original
block--accumulate--accumulate (BAA) codes~\cite{KolesnikovEtAl2026BA}.
Ordinary encoding has similar cost.
Our same-host encoder comparison also includes newer chosen-block BAA,
whose fastest measured configuration takes about $24$\,ms.
In libOTe's regular-noise Silent OT path, the sender produces $2^{18}$ hashed
OTs in $2.944$\,ms, or $89.0$ million OTs per second on one core.
This experiment reuses the full precomputed SPIN code and does not refresh
its permutations between batches; such refresh is a heuristic outside our
distance certificates.
It measures local computation with supplied base correlations,
excluding initial setup and network transport; Section~\ref{sec:spin-pcg}
specifies the timing scope.

As application evidence, we substitute SPIN into a Brakedown-style
PCS~\cite{brakedown23} and integrate it into Flock~\cite{flock26}.
Standalone commitment and one opening of $512$\,MiB take $509$\,ms on one
core. Replacing Ligerito~\cite{ligerito25} within Flock improves throughput
by $1.63$--$2.29\times$ at the measured workloads, with larger proofs and
slower verification. Both backends receive the applicable shared
optimizations. Sections~\ref{sec:spin-pcs} and~\ref{sec:spin-flock}
state the security targets, what they exclude, and the comparison scope.

\subsection{Related Work}

\paragraph{Serial Concatenation.}
Turbo codes use randomized interleaving~\cite{berrou93}.
Benedetto et al.\ analyzed serial concatenation with a uniform
interleaver~\cite{benedetto98}. Repeat--accumulate codes and their
enumerators~\cite{divsalar98}, distance analyses of parallel and serial
concatenations~\cite{kahale98}, and near-GV repeat-multiple-accumulate
ensembles~\cite{kliewer08} form the coding-theoretic background.
Bazzi, Mahdian, and Spielman give limitations for depth-two constant-memory
concatenations and positive results with two accumulators~\cite{bms09}.
Block-Accumulate codes use random block-diagonal maps, detailed expected
spectra, and optimized CPU/GPU encoders~\cite{KolesnikovEtAl2026BA}.
Accumulator SPIN truncates BAA to a single accumulator stage.
Chosen-block BAA reuses a spectrum-controlled constituent~\cite{cryptoeprint:2026/1903};
Structured SPIN combines that shared design idea with a structured interleaver
and a different recursive inner.

\paragraph{Correlation Generation.}
PCGs and VOLE use linear codes
\cite{bcgi18,bcg19pcg,bcg19tworound,ferret20,bcg20ring}.
Silver uses structured LDPC matrices and a banded triangular
solve~\cite{CouteauRindalRaghuraman2021Silver}; its heuristic
linear-distance conjecture was later refuted~\cite{RaghuramanRindalTanguy2023EC}.
Expand--Accumulate combines a sparse expander and accumulator~\cite{ea22}.
Expand--Convolute uses random time-varying convolution and proves distance
through coupling and concentration~\cite{RaghuramanRindalTanguy2023EC};
its faster empirical parameter choices go beyond those guarantees.
Random SPIN uses that inner but replaces the expander with independent
random-block injections and uses an exact weight-enumerator transform.
Structured SPIN instead combines a spectrum-controlled outer, structured
interleaving, and a recursive inner to prove strong distance with linear encoding work.

\paragraph{Proof Systems.}
Linear encoding is a first-order cost in constant-overhead arguments
\cite{rr25}. Orion uses code switching~\cite{orion22}; Brakedown gives
a field-agnostic SNARK from linear-time encodable codes~\cite{brakedown23}.
Other constructions use Expand--Accumulate~\cite{fieldEA24} or
interleaved RAA with improved binary distance bounds~\cite{blaze25}.
SPIN could also instantiate Blaze's generic code-switching framework, which
gives asymptotically smaller proofs than our Brakedown-style PCS.
Bolt uses sketched codes to reduce commitment cost~\cite{bolt26}.

\paragraph{Organization and Artifact.}
Sections~\ref{sec:preliminaries}--\ref{sec:framework} give the common framework;
Section~\ref{sec:spin-intuition} builds intuition before the structured
construction, asymptotic proof, and finite certificates in
Sections~\ref{sec:structured-spin}--\ref{sec:finite-certificates}.
Sections~\ref{sec:implementation} and~\ref{sec:applications} evaluate the encoder
and applications. The appendices contain the complete supporting proofs,
parameter studies, and detailed experiments.
\ifsubmission
The artifact provides the core SPIN encoder implementation and build instructions.
\else
The source repository~\cite{spinCodeRepository} includes the core SPIN
encoder implementation and build instructions.
\fi
The accompanying Lean development includes proof sources, numerical
certificates, and reproduction instructions.
